\documentclass[12pt, a4paper]{article}
\usepackage[czech]{babel}
\usepackage[utf8]{inputenc}
\usepackage[T1]{fontenc}
\usepackage{geometry}
\usepackage{titlesec}
\usepackage{enumitem}
% No global description settings - will use manual formatting
\usepackage[dvipsnames]{xcolor}
\usepackage{xcolor}
\usepackage{hyperref}
\usepackage{fancyhdr}
\usepackage{setspace}

\geometry{left=2.5cm, right=2.5cm, top=2.5cm, bottom=2.5cm}
\onehalfspacing

% Nastavení barev
\definecolor{sectioncolor}{RGB}{0, 102, 204}
\definecolor{subsectioncolor}{RGB}{0, 153, 153}
\definecolor{titlecolor}{RGB}{139, 0, 0}

% Nastavení nadpisů
\titleformat{\section}{\color{sectioncolor}\Large\bfseries}{\thesection}{1em}{}
\titleformat{\subsection}{\color{subsectioncolor}\large\bfseries}{\thesubsection}{1em}{}

% Nastavení stránek
\pagestyle{fancy}
\fancyhf{}
\fancyhead[R]{\color{sectioncolor}\thepage}
\fancyfoot[C]{\color{subsectioncolor}Literární díla - maturitní otázky}
\renewcommand{\headrule}{\hbox to\headwidth{\color{sectioncolor}\leaders\hrule height \headrulewidth\hfill}}
\renewcommand{\footrule}{\hbox to\headwidth{\color{subsectioncolor}\leaders\hrule height \footrulewidth\hfill}}

% Hypertextové odkazy
\hypersetup{
    colorlinks=true,
    linkcolor=sectioncolor,
    urlcolor=subsectioncolor,
    citecolor=titlecolor
}

\newcommand{\mybf}[1]{\textcolor[RGB]{0, 179, 219}{\textbf{#1}}}

\begin{document}

% Titulní strana
\begin{titlepage}
	\centering
	\vspace*{2cm}
	{\color{titlecolor}\Huge\bfseries Literární díla\par}
	\vspace{1cm}
	{\color{sectioncolor}\Large maturitní otázky z českého jazyka a literatury\par}
	\vspace{3cm}
	\vspace{0.5cm}
	\vfill
	{\color{subsectioncolor}\today\par}
\end{titlepage}

\tableofcontents
\newpage

% ============================================================================
\section{Epos o Gilgamešovi}
% ============================================================================

\subsection*{Literární klasifikace}
\mybf{Literární žánr}: Epos - rozsáhlá epická veršovaná skladba\\[3pt]
\mybf{Literární druh}: Epika, poezie, volný verš\\[3pt]
\mybf{Literární forma:} Poezie, text je psaný ve verších. Původně klínovým písmem na hliněných tabulkách v Mezopotámii.\\[10pt]
\mybf{Myšlenka díla}: Smíření se smrtelností, hledání smyslu života, hodnota přátelství, lidská omezenost vůči bohům.\\[10pt]
\mybf{Motiv}: Motiv hledání nesmrtelnosti -- cesta za Uta-napištimem, rostlina mládí, život a osudy vládce Gilgameše, bohové a polobozi, nadpřirozeno, symbolika, strach ze smrti, přátelství, náboženské rysy, odpor vůči bohům, uvědomění -- člověk nemůže být nesmrtelný, nesmrtelný zůstává jen svými činy.\\[10pt]
\mybf{Inspirace}: Gilgameš pravděpodobně skutečně existoval -- byl králem města Uruk v Mezopotámii (cca 2700 př. n. l.). Z reálného panovníka se postupně stal legendární polobůh. Jedná se o reálné dějiny přetavené do mýtu. Sumerská mytologie, starší samostatné mezopotámské legendy.\\[10pt]
\mybf{Kompozice}: Chronologická, dohromady 12 hliněných tabulek, na nichž bylo dílo poprvé nalezeno. Epizodická výstavba -- skládá se z jednotlivých dobrodružných epizod (boj s Chumbabou, Býk nebeský, cesta k Uta-napištimu...).

\subsection*{Vnitřní stavba děje}
\noindent\textbf{Expozice} -- Představení Gilgameše -- mocný, ale krutý král Uruku. Obyvatelé si stěžují bohům. Bohové stvoří Enkidua jako protiváhu.
\noindent\textbf{Kolize} -- Setkání a souboj Gilgameše s Enkiduem. Místo nepřátelství vzniká přátelství. Společně se vydávají na výpravy (Chumbaba, Býk nebeský). Rozhněvají bohy.
\noindent\textbf{Krize} -- Bohové potrestají Enkidua nemocí. Enkidu umírá. Gilgameš je zdrcen a poprvé si plně uvědomuje vlastní smrtelnost.
\noindent\textbf{Peripetie (obrat)} -- Gilgameš se vydává hledat nesmrtelnost. Setkává se s Uta-napištimem (přeživším potopy). Získá rostlinu mládí.
\noindent\textbf{Katastrofa / Rozuzlení} -- Had mu rostlinu ukradne. Gilgameš pochopí, že fyzická nesmrtelnost mu není souzena. Vrací se do Uruku a přijímá svou lidskost.\\[10pt]
\mybf{Doba a místo}: Období raných městských států v Mezopotámii. Konkrétně Uruk (ve starověkém Sumeru, dnešní Irák) kolem 2600 př. n. l., kdy měl historický Gilgameš vládnout. Část děje se odehrává i mimo město: cedrový les (boj s Chumbabou), cesta přes hory, „vody smrti", místo, kde žije Uta-napištim.\\[10pt]
\mybf{Vypravěč}: Er-forma, nezaujatě popisuje děj, promluvy: převážně vyprávění, úvaha.\\[10pt]
\mybf{Jazyk}: Dílo je psáno ve spisovném jazyku, vznešený (patetický) styl, archaismy.\\[10pt]

\subsection*{Postavy}
\begin{itemize}[itemsep=0.3em]
	\item \textbf{Gilgameš:} 5. král sumerského města Uruk, nechal postavit vysoké hradby a nutil všechny pracovat, silný panovník -- panoval přísně až tyransky, ale moudře, vysoce postavený muž nesoucí hrdinské rysy (udatnost, odvaha, chrabrost, nezlomnost)
	\item \textbf{Enkidu:} Původně nepřítel Gilgameše, nakonec ale nejlepší přítel, pomocník, napůl zvíře, byl vyslaný bohy, aby zabil Gilgameše
	\item \textbf{Ištar:} Bohyně, zášť, pýcha, pomstychtivost, chtěla si vzít Gilgameše -- ten odmítl
	\item \textbf{Aruru:} Bohyně, stvořila z hlíny Enkidua
	\item \textbf{Ninsun:} Matka Gilgameše, bohyně -- proto Gilgameš z bůh
	\item \textbf{Chumbaba:} Nestvůra, nadpřirozená bytost, pán cedrového lesa
	\item \textbf{Uta-napištim:} Bůh, dříve člověk, stal se bohem po potopě světa. Ostatní postavy mají roli rádců, nebo naopak protivníků
\end{itemize}

\subsection*{Autor díla}
Vznik z lidové slovesnosti -- autor neznámý. Lidová slovesnost -- příběh si předávalo po staletí několik kultur a proto vzniklo i několik verzí. Dílo bylo nalezeno na hliněných tabulkách na území dnešního Iráku. Zde se nacházelo město Uruk (prvně sumerské, pak babylonské).

\subsection*{Literárně historický kontext}
Nejstarší dílo světové literatury, počátek starověku, vzniká z lidové slovesnosti kolem roku 2000--1700 př. n. l., památka sumerského písemnictví, u nás v té době doba bronzová, nedochován v úplnosti, rekonstruován z více verzí, v této době vznikaly žánry -- mýty, hymny, modlitby, zpěvy, nářky.

{\color{subsectioncolor}\hrulefill}
\subsection*{\color{subsectioncolor}Děj}
\paragraph{1. tabulka}
Když přišel do Uruku, byl Gilgameš krutý vládce („ve dne v noci nutí i děti k práci") a jeho poddaní se obrátili na boha Anu. Ten zavolal Aruru, bohyni lidí, aby stvořila někoho, kdo se postaví Gilgamešovi. Aruru stvořila Enkidua, obdařeného silou od boha války, byl chlupatý jako zvíře a živil se trávou s gazelami. Lovec ho uviděl v horách a oznámil to ostatním lidem. Gilgameš poslal nevěstku Šamchatu, aby ho svedla a odcizila mu jeho zvěř, spolu obcovali sedm dní a nocí a zvířata utekla. Šamchata pak vedla Enkidua do Uruku. Enkidu chtěl s Gilgamešem poměřit síly, ale Šamchata mu to vymlouvala. Gilgameš se o jeho příchodu dozvěděl ze snu a těšil se na něj, jako by se do něj zamiloval.

\paragraph{2. tabulka}
Společně dorazili do města a Šamchata dala Enkiduovi jídlo a pití. Očistila ho a proměnila ze zvířete v člověka (poloboha, vzhledem připomíná Gilgameše). Enkidu potkal Gilgameše v ulicích Uruku a zaklesli se do sebe jako zápasníci, bojovali spolu holýma rukama a Gilgameš poté klesl na kolena. Enkidu uznal jeho výjimečnost a smířili se polibkem -- stali se bratry. Gilgameš navrhl, že spolu zabijí Chumbabu, strážce cedrového lesa. Enkidu se bál, že Chumbabu, boha, který les chrání, nepřekonají, ale Gilgameš ho přemluvil. Připravili se u kovářů na zbraně, rádci cestu vymlouvali, ale Gilgameš je neposlechl a vydal se na riskantní výpravu.

\paragraph{3. tabulka}
Před odchodem navštívili chrám Egalmach v Uruku, sídlo bohyně Ninsun, Gilgamešovy matky. Požádali ji, aby se za ně modlila, a ona souhlasila. Aby upevnila jejich přátelský svazek, adoptovala Enkidua za svého syna, čímž se stal Gilgamešovým bratrem.

\paragraph{4. tabulka}
Po třech dnech cesty vypráví Gilgameš svůj sen, v němž se zřítila hora. Enkidu vysvětluje, že hora symbolizuje Chumbabu a že ho porazí. Následující noc měl sen o stepním býkovi, Enkidu mu vysvětluje, že je to bůh Jamas, který jim pomáhá. Ve třetím snu viděl bouři, blesk a oheň, zbyl jen popel; výklad zněl, že společně porazí Chumbabu. Celkem se Gilgamešovi zdálo pět snů, jejich poselství bylo: sám nic nezvládne, ale ve dvojici mají větší sílu.

\paragraph{5. tabulka}
Dorazili k cedrovému lesu a začali kácet stromy. Chumbaba zaslechl kácení a rozhněval se, vynesl je do výšky, aby je shodil. Bůh Samas jim však pomohl a poslal na Chumbabu silné větry ze všech stran, oslepil ho a Gilgameš s Enkiduem ho zabili třemi ranami. Pokračovali v kácení stromů a posílali je po Eufratu do Uruku; Chumbabovi uřízli hlavu.

\paragraph{6. tabulka}
Po bitvě k Gilgamešovi promluvila Ištar, bohyně lásky, a nabídla mu sňatek. Gilgameš ji však vyčetl její minulost (s ostatními muži), což ji rozhněvalo, a vystoupala na nebesa za svým otcem Anu, aby na něj poslal nebeského býka. Ištar začala ničit Uruk a zabíjet lidi, mezi nimi měl být i Enkidu. Ten však zavolal Gilgameše na pomoc a společně býka porazili. Enkidu vyrval býkovi kýtu a hodil ji po Ištar. Nakonec připravili hostinu pro lid Uruku.

\paragraph{7. tabulka}
Enkidu měl v noci sen: Samas a Anu se radili o tom, proč zabili býka a Chumbabu a proč vykáceli cedry, a rozhodli, že oba musí zemřít. Samas se je snažil ochránit a vysvětlil, že vše bylo na jeho rozkaz. Bohové poté ustoupili na to, že zemřít musí pouze Enkidu. Enkidu litoval vykácených stromů a modlil se s Gilgamešem k bohům za odpuštění, ale bohové své rozhodnutí nezměnili. V zármutku proklínal Šamchatu, ale Samas ho kárá a varuje, že nemá zlorečit té, která mu dala jídlo a seznámila ho s Gilgamešem, Enkidu ji tedy požehnal. Poté měl další sen o Anzu, ptáku bouře (orlí tělo, hlava lva), který ho uchopil a hodil do podsvětí, ve snu viděl, kdo v podsvětí všechno je. Enkidu onemocněl, ležel 12 dní a nakonec zemřel, přičemž vinil Gilgameše, že mu nepomohl.

\paragraph{8. tabulka}
Tato tabulka začíná dlouhým žalozpěvem Gilgameše, ve kterém vyjmenovává, co spolu s Enkiduem prožili, a kdo všechno by ho měl oplakávat. Poté proplakal sedm dní a sedm nocí. Nakonec vyzval kováře a zlatníky, aby vyrobili sochu Enkidua, a poté ho konečně pohřbil. Gilgameš se poté ptá boha Samase na smysl života a vydává se na step, kde bloudí. Samas k němu mluví: „Kam běžíš? Život, který hledáš, nenalezneš" -- odkaz na nemožnost získat nesmrtelnost.

\paragraph{9. tabulka}
Gilgameš utíká do stepi a bojí se smrti, modlí se k bohům. Škorpion-člověk se ho ptá, proč zavítal do tak vzdálené stepi. Gilgameš odpovídá, že jde za jedním z bohů zeptat se na smysl života a smrti (Uta-napištim). Škorpion-člověk ho varuje, že ho čeká dlouhá a trnitá cesta přes hory, ale Gilgameš ho přemluví a ten mu otevře cestu. Po trnité stezce a dni v lese bez světla dorazí do krásné zahrady, kde rostou stromy z diamantů.

\paragraph{10. tabulka}
Gilgameš dorazil k domu Šenkýrky, která se ho ptá, co pohledává tak daleko od Uruku. Vypráví jí o svém příteli a o tom, že neví, co dál dělat. Šenkýrka mu radí, aby si užíval radosti života, ale Gilgameš ji neposlouchá a hledá způsob, jak přeplout moře, aby se zeptal Uta-napišti, jediného přeživšího potopy, na nesmrtelnost. Šenkýrka mu poradí, že jediný, kdo ho může přepravit, je lodivod Sursunab. Gilgameš ho najde a vysvětlí mu, že ztratil přítele a nechce dopadnout stejně; Sursunab ho vezme na loď a nařídí mu, aby si vzal dlouhé tyče. Po měsíci dorazí k Vodám smrti, kde Gilgameš používá tyče k odrážení od dna, protože se nesmí dotknout vody. Lodivod mu připomene, že nic na Zemi netrvá věčně a že Gilgameš, byť výjimečný, nebyl stvořen jako bůh a stejně jednou zemře.

\paragraph{11. tabulka}
Gilgameš se ptá Uta-napišti, jediného přeživšího potopy, jak se stal nesmrtelným. Uta-napišti mu vypráví, že bůh sladkých vod Ea ho varoval před potopou a poradil mu postavit archu ve tvaru čtverce, na kterou naložil rodinu, stříbro a stepní zvěř. Bohové spustili bouři a potopu, po sedmi dnech bouře ustala, ale loď zůstala šest dní uvízlá na hoře. Uta-napišti vypustil tři ptáky, kteří se vrátili, jelikož nenašli zem, pak zapálil oběť bohům a Enlil mu dal s manželkou požehnání a učinil je bohy, tím nebyl jediný člověk na zemi a civilizace mohla znovu začít. Gilgameš se mohl podrobit zkoušce nesmrtelnosti -- týden nespat -- ale ihned usnul. Uta-napišti mu poradil, jak se zbavit strachu ze smrti: lodivod Sursunab ho měl očistit, převléct a vrátit do Uruku. Poté mu Uta-napišti řekl o rostlině a Gilgameš se potopil pro rostlinu vodní růže, která měla přinést radost života. Na cestě zpátky do Uruku mu ji ovšem sežral had a proměnil se ve lva. Po návratu do Uruku bez rostliny si však Gilgameš naplno uvědomil, že jeho nesmrtelnost spočívá v jeho činech, když přehlédl jeho hradby, které postavil a chrání tak své lidi před nepřáteli.

\paragraph{12. tabulka}
Dvanáctá tabulka je velmi poškozená a fragmentární, většina textu se ztratila, takže obsahuje jen zlomky a není kompletní. Překladatelé často dávají jen poznámku o její existenci a pár dochovaných útržků. Předpokládá se, že se v ní Gilgameš vrací domů do Uruku a vypravuje o svém hledání nesmrtelnosti. Není jasné, jestli v originále objevoval nějaký „morální závěr" nebo moudrost, kterou získal -- fragmenty naznačují jen, že jeho poznání bylo částečně ztraceno.
\newpage

% ============================================================================
\section{Romeo a Julie}
% ============================================================================

\subsection*{Literární klasifikace}
\mybf{Literární žánr:} Tragédie (milostná, veršovaná, o 5 dějstvích)\\[3pt]
\mybf{Literární druh:} Drama\\[3pt]
\mybf{Funkční styl:} Umělecký\\[3pt]
\mybf{Slohový postup:} Vyprávěcí\\[10pt]
\mybf{Myšlenka díla:} Autor poukazuje na to, aby mladí lidé, stejně jako Romeo a Julie, bojovali za své právo na lásku, štěstí a právo rozhodovat o vlastním osudu i přes jeho nepřízeň. Když se nepřátelé nechtějí usmířit, odnesou to často nevinní.\\[10pt]
\mybf{Motiv:} Láska, souboje, smrt, vdavky, válka rodů, dlouhý spor.\\[10pt]
\mybf{Inspirace}: Inspirací pro dílo byla středověká a renesanční literatura o tragické lásce, zejména italské pověsti. Hlavním zdrojem byla básnický příběh \textit{The Tragical History of Romeus and Juliet} od Arthura Brookea z roku 1562. Shakespeare příběh přetvořil do dramatické podoby, rozvinul postavy, dialogy a situace tak, aby vznikla tragédie nadčasová a univerzální, která zobrazuje sílu lásky proti nenávisti a rodinným konfliktům.\\[10pt]
\mybf{Kompozice:} Chronologická kompozice, člení se na prolog a 5 dějství.

\subsection*{Vnitřní stavba děje}
\noindent\textbf{Expozice:} Seznamujeme se s prostředím a konfliktem. Ve Veroně se dlouhodobě nenávidí rody Monteků a Kapuletů. Poznáváme hlavní postavy. Romeo a Julie se poprvé potkají na plese u Kapuletů a zamilují se.
\noindent\textbf{Kolize:} Konflikt se začíná rozvíjet. Romeo a Julie zjistí, že patří k nepřátelským rodům. Přesto se rozhodnou tajně vzít.
\noindent\textbf{Krize:} Situace se vyhrocuje a komplikuje. Tybalt zabije Mercutia. Romeo v pomstě zabije Tybalta a je vyhnán z Verony. Julie má být provdána za Parise.
\noindent\textbf{Peripetie (zvrat):} Dochází k zásadnímu obratu v ději. Julie vypije lektvar od mnicha Vavřince, aby předstírala smrt. Plán je, že Romeo ji potom probudí. Romeovi se ale zpráva nedostane.
\noindent\textbf{Katastrofa:} Tragický konec. Romeo si myslí, že Julie je opravdu mrtvá, a vypije jed. Julie se probudí a spáchá sebevraždu dýkou. Smrt milenců nakonec usmíří oba rody.\\[10pt]

\mybf{Doba a místo:} Italské město Verona ve druhé polovině 16. století, a později i Mantova.\\[10pt]
\mybf{Vypravěč:} Jelikož se jedná o literární druh drama, nevyskytuje se v tomto díle vypravěč. Příběh tvoří promluvy jednotlivých postav a důležité vševědoucí informace doplňují scénické poznámky.\\[10pt]
\mybf{Jazyk:} Dílo je psáno formou dramatu, proto je děj tvořen především dialogy a monology postav. Text je napsán pětistopým jambickým veršem, tzv. blankversem, který je pro Shakespearova díla typický. V jazyce se objevuje spisovná a místy knižní slovní zásoba, archaismy a básnické prostředky, například metafory, personifikace a přirovnání. Časté jsou tázací a zvolací věty, které zvyšují dramatičnost textu. Vyskytuje se také inverze, tedy nezvyklé pořadí slov ve větě.

\subsection*{Postavy}
\begin{itemize}[itemsep=0.3em]
	\item \textbf{Romeo:} Mladý, divoký, milý, zamiluje se bezhlavě do Julie, ale předtím už byl jednou také hodně zamilovaný, jeho láska možná vůbec není stálá. Vývoj -- z nezodpovědného floutka opravdový muž, chtěl, aby boje mezi rody ustaly.
	\item \textbf{Montek a Monteková:} Romeovi rodiče.
	\item \textbf{Benvolio:} Montekův synovec a Romeův přítel.
	\item \textbf{Julie:} Velice mladinká, krásná dívka, rodiče jí našli ženicha, za kterého se má brzy provdat, ona se však zamiluje do Romea. Vývoj -- z mladé nezkušené dívky se stává dospělá žena, která je schopna pro svou lásku obětovat cokoliv.
	\item \textbf{Kapulet:} Juliin otec, v domě vládne tvrdou rukou, chce, aby ho Julie jako správná dcera ve všem na slovo poslouchala; pokud ne, je schopen se hrozně rozzuřit a dceru vykázat z domu.
	\item \textbf{Kapuletová:} Juliina matka, chová se ke své dceři velice odměřeně, až chladně.
	\item \textbf{Tybalt:} Juliin bratranec, libuje si v bitkách mezi Kapulety a Monteky.
	\item \textbf{Kníže Eskalus:} Veronský vládce, už mu dochází trpělivost se znepřátelenými rody.
	\item \textbf{Chůva:} Opatruje a vychovává Julii od jejího narození, má ji ráda jako svou dceru, udělala by pro ni cokoliv.
	\item \textbf{Merkucio:} Romeův nejlepší přítel, je starší -- dohlíží na Romea, je mu vzorem.
	\item \textbf{Paris:} Mladý hrabě, příbuzný knížete.
	\item \textbf{Bratr Lorenzo a bratr Baltazar:} Františkáni.
	\item \textbf{Baltazar:} Romeův sluha.
	\item \textbf{Samson a Řehoř:} Kapuletovi sluhové.
	\item \textbf{Petr:} Sluha Juliiny chůvy.
	\item \textbf{Abraham:} Montekův sluha.
	\item \textbf{Další postavy:} Lékárník, stařec, tři muzikanti, veronští občané, členové Kapuletovy a Montekovy rodiny, maskovaní hosté, stráže, družina, Chorus.
\end{itemize}

\subsection*{Autor díla}
\textbf{William Shakespeare} byl anglický básník a dramatik, narozený ve Stratfordu nad Avonou. Za divadelní kariérou odchází do Londýna, kde vyniká jako herec a dramatik a stává se spolumajitelem divadla Globe, kruhového divadla bez opon a kulis s blízkým kontaktem s diváky. Závěr života strávil opět ve Stratfordu. Jeho tvorba se dělí na:
\begin{itemize}[itemsep=0.3em]
	\item \textbf{Komedie:} \textit{Zkrocení zlé ženy}, \textit{Kupec benátský}, \textit{Mnoho povyku pro nic}
	\item \textbf{Historické hry:} \textit{Král Richard III.}, \textit{Král Jindřich IV.}
	\item \textbf{Tragédie:} \textit{Romeo a Julie}, \textit{Macbeth}, \textit{Hamlet}, \textit{Othello} -- považovány za vrchol anglické literatury
\end{itemize}
Díla jsou přelomová díky propracované psychologii postav a realistickému pohledu na osud -- postavy si určují svůj vlastní osud a tragédie vzniká z lidské vášně, nikoli z náhody nebo boží vůle. Jazyk střídá verš s prózou a typickým prostředkem je blankvers, nerymovaný pětistopý jambický verš s převahou mužského zakončení. Shakespeare napsal 37 dramat a sbírku sonetů.

\subsection*{Literárně historický kontext}
V tomto období se objevují dva významné směry:
\noindent\textbf{Renesance} -- Znamená „znovuzrození" a obnovení zájmu o antiku. Měšťané se odvracejí od středověkého způsobu myšlení a odříkavého životního stylu a přiklání se k radostem pozemského života. Prosazuje se rozvoj přírodních věd a věd o člověku, středem zájmu je člověk, nikoli Bůh. Rozvíjí se stavitelství a umění -- stavěly se zámky a paláce, malíři zachycovali stav lidské duše a ženskou krásu. Do literatury pronikají nové literární žánry a národní jazyky.
\noindent\textbf{Humanismus} -- Zdůrazňuje zájem o člověka, víru v jeho schopnosti, vlastnosti a rozum, který umožňuje poznávat svět. Klade důraz na svobodný vývoj člověka.
Tyto dva směry mají mnoho společného; renesance je více zaměřená na pozemské radosti a obsahuje více motivů lásky, erotiky, přátelství a podobně.
{\color{subsectioncolor}\hrulefill}
\subsection*{\color{subsectioncolor}Děj}
Děj se odehrává v severní Itálii, ve městě Veroně. Zde žily dva znepřátelené rody -- Kapuletové, kteří měli třináctiletou dceru Julii, a Montekové, kteří měli syna Romea. Jednoho dne sluhové obou rodů vyvolají další rozepři. Vévoda veronský jim kvůli tomu uloží ultimátum. Ještě jedna rozepře a zaplatí smrtí. Kapuletovi plánují zásnuby své dcery s mladým šlechticem Parisem.

Romeo je nešťastně zamilován do Rosalie, a proto ho pro rozptýlení vezmou jeho přátelé na ples Kapuletů. Zde spatří Julii a okamžitě se do sebe oba zamilují. Během noci se rozhodnou, že se vezmou, a tak jsou bratrem Vavřincem tajně oddáni. Druhý den dojde k další rozepři, kde Tybalt zabije Merkucia, Romeova přítele. Později Tybalt vyprovokuje k boji i Romea, ten pomstí smrt svého přítele a Tybalta zabije. Romeo musí odejít z Verony, protože porušil dohodu s vévodou.

Mezitím se plánuje svatba Julie a Parise. Julie však miluje Romea, nechce si Parise vzít, proto jde pro pomoc za Vavřincem. Ten jí dá nápoj, který Julii uspí a po jehož požití bude dva dny vypadat jako mrtvá. V den svatby s Parisem vypije Julie lektvar. Vypadá jako mrtvá, Kapuletové i Paris truchlí a pohřbí ji v rodinné hrobce. Vavřinec na nic nečeká a posílá ihned zprávu Romeovi. Než však sluha stačí zprávu vyřídit, dozví se Romeo o Juliině smrti od sluhy Baltazara, který neví nic o tajném plánu.

Romeo okamžitě jede k hrobce Kapuletů, kde se setká s Parisem, kterého tam po hádce zabije. Když Romeo spatří v hrobce mrtvou Julii, touží zemřít po jejím boku. Tak vypije velmi silný jed a umírá. V té chvíli se probouzí Julie a jak spatří mrtvého Romea, probodne se jeho dýkou. Když k hrobce přijdou Kapuletové a Montek, jehož žena zemřela, když se dozvěděla o smrti syna, začne jim Vavřinec vysvětlovat celý příběh. Oba rody si uvědomují, že kvůli svým sporům přišly o své děti a tak si nad jejich mrtvými těly podají ruce, čímž mezi nimi opět zavládne mír.

\newpage


% ============================================================================
\section{Lakomec}
% ============================================================================

\subsection*{Literární klasifikace}
\mybf{Literární druh:} \\[3pt]
\mybf{Literární forma:} \\[3pt]
\mybf{Funkční styl:} \\[3pt]
\mybf{Slohový postup:} \\[10pt]
\mybf{Myšlenka díla:} \\[10pt]
\mybf{Motiv:} \\[10pt]
\mybf{Inspirace:} \\[10pt]
\mybf{Kompozice:} 

\subsection*{Vnitřní stavba děje}
\noindent\textbf{Expozice} 
\noindent\textbf{Kolize} 
\noindent\textbf{Krize} 
\noindent\textbf{Peripetie (zvrat)} 
\noindent\textbf{Katastrofa} \\[10pt]
\mybf{Doba a místo:} \\[10pt]
\mybf{Vypravěč:} \\[10pt]
\mybf{Jazyk:} 

\subsection*{Postavy}
\begin{itemize}[itemsep=0.3em]
	\item \textbf{Napoleon:}
\end{itemize}

\subsection*{Autor díla}
\textbf{Eric Arthur Blair (1903--1950)}

\textbf{Další díla:} \textit{Barmské dny}, \textit{Farářova dcera}, \textit{1984}, \textit{Na dně v Paříži a Londýně}, \textit{Nadechnout se}.

\subsection*{Literárně historický kontext}

{\color{subsectioncolor}\hrulefill}
\subsection*{\color{subsectioncolor}Děj}

\newpage


% ============================================================================
\section{Revizor}
% ============================================================================

\subsection*{Literární klasifikace}
\mybf{Literární druh:} \\[3pt]
\mybf{Literární forma:} \\[3pt]
\mybf{Funkční styl:} \\[3pt]
\mybf{Slohový postup:} \\[10pt]
\mybf{Myšlenka díla:} \\[10pt]
\mybf{Motiv:} \\[10pt]
\mybf{Inspirace:} \\[10pt]
\mybf{Kompozice:} 

\subsection*{Vnitřní stavba děje}
\noindent\textbf{Expozice} 
\noindent\textbf{Kolize} 
\noindent\textbf{Krize} 
\noindent\textbf{Peripetie (zvrat)} 
\noindent\textbf{Katastrofa} \\[10pt]
\mybf{Doba a místo:} \\[10pt]
\mybf{Vypravěč:} \\[10pt]
\mybf{Jazyk:} 

\subsection*{Postavy}
\begin{itemize}[itemsep=0.3em]
	\item \textbf{Napoleon:}
\end{itemize}

\subsection*{Autor díla}
\textbf{Eric Arthur Blair (1903--1950)}

\textbf{Další díla:} \textit{Barmské dny}, \textit{Farářova dcera}, \textit{1984}, \textit{Na dně v Paříži a Londýně}, \textit{Nadechnout se}.

\subsection*{Literárně historický kontext}

{\color{subsectioncolor}\hrulefill}
\subsection*{\color{subsectioncolor}Děj}

\newpage


% ============================================================================
\section{Král Lávra}
% ============================================================================

\subsection*{Literární klasifikace}
\mybf{Literární druh:} \\[3pt]
\mybf{Literární forma:} \\[3pt]
\mybf{Funkční styl:} \\[3pt]
\mybf{Slohový postup:} \\[10pt]
\mybf{Myšlenka díla:} \\[10pt]
\mybf{Motiv:} \\[10pt]
\mybf{Inspirace:} \\[10pt]
\mybf{Kompozice:} 

\subsection*{Vnitřní stavba děje}
\noindent\textbf{Expozice} 
\noindent\textbf{Kolize} 
\noindent\textbf{Krize} 
\noindent\textbf{Peripetie (zvrat)} 
\noindent\textbf{Katastrofa} \\[10pt]
\mybf{Doba a místo:} \\[10pt]
\mybf{Vypravěč:} \\[10pt]
\mybf{Jazyk:} 

\subsection*{Postavy}
\begin{itemize}[itemsep=0.3em]
	\item \textbf{Napoleon:}
\end{itemize}

\subsection*{Autor díla}
\textbf{Eric Arthur Blair (1903--1950)}

\textbf{Další díla:} \textit{Barmské dny}, \textit{Farářova dcera}, \textit{1984}, \textit{Na dně v Paříži a Londýně}, \textit{Nadechnout se}.

\subsection*{Literárně historický kontext}

{\color{subsectioncolor}\hrulefill}
\subsection*{\color{subsectioncolor}Děj}

\newpage


% ============================================================================
\section{Kytice}
% ============================================================================

\subsection*{Literární klasifikace}
\mybf{Literární druh:} \\[3pt]
\mybf{Literární forma:} \\[3pt]
\mybf{Funkční styl:} \\[3pt]
\mybf{Slohový postup:} \\[10pt]
\mybf{Myšlenka díla:} \\[10pt]
\mybf{Motiv:} \\[10pt]
\mybf{Inspirace:} \\[10pt]
\mybf{Kompozice:} 

\subsection*{Vnitřní stavba děje}
\noindent\textbf{Expozice} 
\noindent\textbf{Kolize} 
\noindent\textbf{Krize} 
\noindent\textbf{Peripetie (zvrat)} 
\noindent\textbf{Katastrofa} \\[10pt]
\mybf{Doba a místo:} \\[10pt]
\mybf{Vypravěč:} \\[10pt]
\mybf{Jazyk:} 

\subsection*{Postavy}
\begin{itemize}[itemsep=0.3em]
	\item \textbf{Napoleon:}
\end{itemize}

\subsection*{Autor díla}
\textbf{Eric Arthur Blair (1903--1950)}

\textbf{Další díla:} \textit{Barmské dny}, \textit{Farářova dcera}, \textit{1984}, \textit{Na dně v Paříži a Londýně}, \textit{Nadechnout se}.

\subsection*{Literárně historický kontext}

{\color{subsectioncolor}\hrulefill}
\subsection*{\color{subsectioncolor}Děj}

\newpage

% ============================================================================
\section{Divá Bára}
% ============================================================================

\subsection*{Literární klasifikace}
\mybf{Literární druh:} \\[3pt]
\mybf{Literární forma:} \\[3pt]
\mybf{Funkční styl:} \\[3pt]
\mybf{Slohový postup:} \\[10pt]
\mybf{Myšlenka díla:} \\[10pt]
\mybf{Motiv:} \\[10pt]
\mybf{Inspirace:} \\[10pt]
\mybf{Kompozice:} 

\subsection*{Vnitřní stavba děje}
\noindent\textbf{Expozice} 
\noindent\textbf{Kolize} 
\noindent\textbf{Krize} 
\noindent\textbf{Peripetie (zvrat)} 
\noindent\textbf{Katastrofa} \\[10pt]
\mybf{Doba a místo:} \\[10pt]
\mybf{Vypravěč:} \\[10pt]
\mybf{Jazyk:} 

\subsection*{Postavy}
\begin{itemize}[itemsep=0.3em]
	\item \textbf{Napoleon:}
\end{itemize}

\subsection*{Autor díla}
\textbf{Eric Arthur Blair (1903--1950)}

\textbf{Další díla:} \textit{Barmské dny}, \textit{Farářova dcera}, \textit{1984}, \textit{Na dně v Paříži a Londýně}, \textit{Nadechnout se}.

\subsection*{Literárně historický kontext}

{\color{subsectioncolor}\hrulefill}
\subsection*{\color{subsectioncolor}Děj}

\newpage

% ============================================================================
\section{Maryša}
% ============================================================================

\subsection*{Literární klasifikace}
\mybf{Literární žánr}: Tragédie\\[3pt]
\mybf{Literární druh}: Drama\\[3pt]
\mybf{Literární forma}: Dialogické drama -- text tvořený především dialogy mezi postavami, s krátkými scénickými poznámkami.\\[3pt]
\mybf{Funkční styl}: Umělecký (literární)\\[3pt]
\mybf{Slohový postup}: Vyprávění a popis\\[10pt]
\mybf{Myšlenka díla}: Moc peněz, síla předsudků, vynucená pokora a poslušnost i ztráta možnosti svobodně nakládat se svým životem přináší jen neštěstí. Vůle člověka by měla být respektována.\\[10pt]
\mybf{Motiv}: Manželství z donucení, útlak ženy, osud vs. svobodná vůle, tragédie jednotlivce, rodinné a společenské tlaky.\\[10pt]
\mybf{Inspirace}: Realistická venkovská společnost 19. století, zkušenosti autorů s moravským venkovem, sociální problémy a omezení žen v tehdejší době.\\[10pt]
\mybf{Kompozice}: Chronologická
\subsection*{Vnitřní stavba děje}: Děj se dělí na 5 jednání:
\vspace{0.5em}
\noindent\textbf{1. jednání -- expozice:} Seznámení s postavami a dějem.
\noindent\textbf{2. + 3. jednání -- kolize:} Maryša svoluje k sňatku, Francek se vrací a vrcholí krize.
\noindent\textbf{4. jednání -- peripetie:} Pokus zabít Francka.
\noindent\textbf{5. jednání -- katastrofa:} Otrávení Vávry.\\[10pt]
\mybf{Doba a místo}: 19. století, moravská vesnice.\\[10pt]
\mybf{Vypravěč}: Neexistuje -- příběh je vyprávěn skrze dialogy a scénické poznámky.\\[10pt]
\mybf{Jazyk}: Autoři používali nářečí z jižní Moravy, slovácké nářečí z okolí slovenských hranic. Používají stylizované moravské nářečí, kombinující prvky hanáckého a moravsko-slovanského nářečí, zejména nářečí Brněnska a konkrétně slovácké nářečí.\\[10pt]
\subsection*{Postavy}
\begin{itemize}[itemsep=0.3em]
	\item \textbf{Maryša:} Lízalova dvacetiletá dcera, později Vávrova manželka; milující žena, která trpí pro svou lásku.
	\item \textbf{Lízal:} Sedlák, který dříve více dbal na peníze než na štěstí svého dítěte.
	\item \textbf{Lízalka:} Maryšina matka; žena s kamennou tváří, která nezná lítost ani pochopení vůči dceři.
	\item \textbf{Vávra Filip:} Muž, který nehledí na štěstí své ženy, ale na její věno; vdovec se třemi dětmi.
	\item \textbf{Francek (František):} Muž, kterého Maryša miluje; říká jen to, co si opravdu myslí.
\end{itemize}

\subsection*{Autoři díla}
\textbf{Alois Mrštík (1861--1925)} -- Český prozaik a dramatik. Vyučoval v několika obcích na jižní Moravě. Přispíval do časopisu Máj, Lumír, Světozor atd.

\textbf{Další tvorba:} sbírka povídek \textit{Nit stříbrná}, povídky \textit{Dobré duše}.

\textbf{Vilém Mrštík (1863--1912)} -- Český překladatel, kritik, prozaik, dramatik. Své studium nedokončil. Věnoval se publicistice. Vyzdvihuje mravní hodnoty.

\textbf{Další tvorba:} román \textit{Santa Lucia}, literární kritika \textit{Moje sny}.

\subsection*{Literárně historický kontext}
Česká realistická literatura 19. století; reflektuje venkovský život, patriarchální společnost a omezení žen; dílo je součástí realismu zaměřeného na sociální kritiku.

{\color{subsectioncolor}\hrulefill}
\subsection*{\color{subsectioncolor}Děj}
Sedlák Lízal, který myslí hlavně na peníze, chtěl provdat svou dceru za mlynáře Vávru. Maryša je ale zamilovaná do prostého chlapce Francka. Vávra slibuje, že se o Maryšu dobře postará, doopravdy mu však jde jen o její bohaté věno, ze kterého chce zaplatit své dluhy. Po svatbě se začne Vávra opíjet, nestará se o rodinu a soudí se s Lízalem o peníze z věna. Sedlákovi konečně dojde, jak nešťastně provdal svou dceru a peníze mu odmítne dát. Francek začne plánovat útěk s Maryšou do Brna, kde našel pro oba práci. Maryša nechce odejít. Opilý Vávra se pokusí Francka zabít, ale Maryša ho zachrání. Maryša je již zoufalá a nasype Vávrovi do kávy jed, pak se k tomu sama přizná.

\newpage

% ============================================================================
\section{Pes baskervillský}
% ============================================================================

\subsection*{Literární klasifikace}
\mybf{Literární druh:} \\[3pt]
\mybf{Literární forma:} \\[3pt]
\mybf{Funkční styl:} \\[3pt]
\mybf{Slohový postup:} \\[10pt]
\mybf{Myšlenka díla:} \\[10pt]
\mybf{Motiv:} \\[10pt]
\mybf{Inspirace:} \\[10pt]
\mybf{Kompozice:} 

\subsection*{Vnitřní stavba děje}
\noindent\textbf{Expozice} 
\noindent\textbf{Kolize} 
\noindent\textbf{Krize} 
\noindent\textbf{Peripetie (zvrat)} 
\noindent\textbf{Katastrofa} \\[10pt]
\mybf{Doba a místo:} \\[10pt]
\mybf{Vypravěč:} \\[10pt]
\mybf{Jazyk:} 

\subsection*{Postavy}
\begin{itemize}[itemsep=0.3em]
	\item \textbf{Napoleon:}
\end{itemize}

\subsection*{Autor díla}
\textbf{Eric Arthur Blair (1903--1950)}

\textbf{Další díla:} \textit{Barmské dny}, \textit{Farářova dcera}, \textit{1984}, \textit{Na dně v Paříži a Londýně}, \textit{Nadechnout se}.

\subsection*{Literárně historický kontext}

{\color{subsectioncolor}\hrulefill}
\subsection*{\color{subsectioncolor}Děj}

\newpage


% ============================================================================
\section{Proces}
% ============================================================================

\subsection*{Literární klasifikace}
\mybf{Literární žánr:} Existencialistický / prozaický román\\[3pt]
\mybf{Literární druh:} Epika\\[3pt]
\mybf{Literární forma:} Próza\\[3pt]
\mybf{Funkční styl:} Umělecký\\[3pt]
\mybf{Slohový postup:} Vyprávěcí\\[10pt]
\mybf{Myšlenka díla:} Válka člověka se sebou samým a rovněž s byrokracií, nejasný boj se soudními praktikami, reflektuje rozbouřené dvacáté století a hledání smyslu života.\\[10pt]
\mybf{Motiv:} Bezmocnost, zmatení, soud, rozsudek, smrt.\\[10pt]
\mybf{Inspirace:} Inspirací k románu Proces byly hlavně jeho osobní zkušenosti. Kafka pracoval v pojišťovně, kde poznal složitou byrokracii. Ovlivnil ho také jeho přísný otec a nešťastný vztah se snoubenkou Felice Bauer. Tyto zkušenosti se promítly do příběhu o člověku, který je bez vysvětlení souzen.\\[10pt]
\mybf{Kompozice:} Chronologická.

\subsection*{Vnitřní stavba děje}
\vspace{0.5em}
\noindent\textbf{Expozice} -- Hlavní postava Josef K. je ráno bez vysvětlení zatčen, i když může dál chodit do práce. Dozví se, že proti němu probíhá jakýsi zvláštní soudní proces, ale nikdo mu neřekne, z čeho je obviněn.
\noindent\textbf{Kolize (zápletka)} -- Josef K. se snaží zjistit, jaký zločin spáchal a jak se může bránit. Začíná chodit k soudu, který je chaotický, byrokratický a nesrozumitelný.
\noindent\textbf{Peripetie (stupňování děje)} -- Josef K. se setkává s různými lidmi spojenými se soudem (advokát, malíř Titorelli, úředníci). Zjišťuje, že systém je neproniknutelný a absurdní a že jeho šance na osvobození jsou téměř nulové. Postupně ztrácí jistotu i kontrolu nad svým životem.
\noindent\textbf{Krize} -- Josef K. si uvědomuje, že proces nikdy skutečně neskončí a že je do něj úplně uvězněn.
\noindent\textbf{Katastrofa (závěr)} -- Dva muži Josefa K. odvedou na okraj města a bez vysvětlení ho zavraždí nožem.

\subsection*{Autor díla}
\textbf{Franz Kafka (1883--1924)} -- jeden z nejvýznamnějších představitelů moderní literatury. Předznamenal vlnu absurdního dramatu. Narodil se v Praze v bohaté židovské rodině. Těžce nesl svou práci jako úředník v pojišťovně. Měl v sobě rozpor, zda se věnovat rodinnému životu nebo psaní, kterému šlo nejlépe v izolaci. Onemocněl tuberkulózou. Před smrtí požádal svého přítele Maxe Broda, aby spálil všechny jeho rukopisy. Řadí se mezi autory pražské německé literatury. Cítil se podivně izolován jako německy píšící židovský autor v Praze. Nepsal pro čtenáře, ale pro sebe. Jeho hrdinové se dostávají do absurdních situací. Autor nevysvětluje proč, ale vše líčí do detailu. Kořeny můžeme hledat v Kafkově citlivé povaze a křehké psychice.

\textbf{Autorovo další dílo:}
\begin{itemize}[itemsep=0.3em]
	\item \textbf{Romány:} \textit{Nezvěstný}, \textit{Proměna}, \textit{Zámek}
\end{itemize}

\subsection*{Literárně historický kontext}
Napjaté vztahy mezi evropskými velmocemi -- zformování Trojspolku (Německé císařství, Rakousko-Uhersko, Osmanská říše, Bulharsko (Itálie)) a Dohody (Francie, VB, Rusko + později Itálie a USA).

	{\color{subsectioncolor}\hrulefill}
\subsection*{\color{subsectioncolor}Děj}
Román Proces se odehrává v prostředí soudu, které Kafka jako vystudovaný právník dobře znal. Josef K. je v den svých třicátých narozenin zatčen z neznámého důvodu a vše probíhá velmi prapodivně. Bankovní úředník Josef K. o žádném provinění neví, během procesu se snaží obhájit svou nevinu. Postupně získává víc informací a zjišťuje, že proces nelze zrušit. Žádný rozsudek není vynesen, ale den před svými třicátými prvními narozeninami je dvěma muži vyveden za město, kde je zabit ranou nožem do srdce.

\newpage

% ============================================================================
\section{Proměna}
% ============================================================================

\subsection*{Literární klasifikace}
\mybf{Literární druh:} \\[3pt]
\mybf{Literární forma:} \\[3pt]
\mybf{Funkční styl:} \\[3pt]
\mybf{Slohový postup:} \\[10pt]
\mybf{Myšlenka díla:} \\[10pt]
\mybf{Motiv:} \\[10pt]
\mybf{Inspirace:} \\[10pt]
\mybf{Kompozice:} 

\subsection*{Vnitřní stavba děje}
\noindent\textbf{Expozice} 
\noindent\textbf{Kolize} 
\noindent\textbf{Krize} 
\noindent\textbf{Peripetie (zvrat)} 
\noindent\textbf{Katastrofa} \\[10pt]
\mybf{Doba a místo:} \\[10pt]
\mybf{Vypravěč:} \\[10pt]
\mybf{Jazyk:} 

\subsection*{Postavy}
\begin{itemize}[itemsep=0.3em]
	\item \textbf{Napoleon:}
\end{itemize}

\subsection*{Autor díla}
\textbf{Eric Arthur Blair (1903--1950)}

\textbf{Další díla:} \textit{Barmské dny}, \textit{Farářova dcera}, \textit{1984}, \textit{Na dně v Paříži a Londýně}, \textit{Nadechnout se}.

\subsection*{Literárně historický kontext}

{\color{subsectioncolor}\hrulefill}
\subsection*{\color{subsectioncolor}Děj}

\newpage

% ============================================================================
\section{Stařec a moře}
% ============================================================================

\subsection*{Literární klasifikace}
\mybf{Literární druh:} Epika\\[3pt]
\mybf{Literární žánr:} Novela\\[3pt]
\mybf{Literární forma:} Próza\\[3pt]
\mybf{Myšlenka díla:} nezlomnost lidského ducha, lidská důstojnost a vnitřní síla. Hemingway ukazuje, že i když člověk prohraje materiálně (Santiago přijde o celou rybu), může vyjít ze souboje jako morální vítěz.\\[10pt]
\mybf{Motiv:} Jeden z nejvýraznějších a nejkrásnějších motivů novely jsou lvi hrající si na africkém pobřeží.Santiago o těchto lvech často sní, když po těžké práci usne. Jde o vzpomínku na doby jeho mládí, kdy se jako námořník plavil k břehům Afriky a pozoroval lvíčata hrající si v podvečer na pláži.\\[10pt]
\mybf{Inspirace:} Reálná událost z roku 1936: Hemingway se o podobném případu dozvěděl už 16 let před vydáním knihy. V roce 1936 publikoval v magazínu Esquire esej s názvem Na modré vodě (Dopis z Golfského proudu). Popsal v ní historku o starém kubánském rybáři, který bojoval obrovského marlína, zaneslo ho to na širé moře a jeho úlovek mu sežrali žraloci. Rybáře pak našli napůl šíleného plakat ve člunu.\\[10pt]
\mybf{Kompozice:} Nemá kapitoly, chronoligický postup

\subsection*{Vnitřní stavba děje}
\noindent\textbf{Expozice:} Seznámení se starým kubánským rybářem Santiagem, který má obrovskou smůlu – už 84 dní neulovil jedinou rybu (je takzvaně „salao“). Zjišťujeme, v jak chudých podmínkách žije a jaký má vztah k chlapci Manolinovi, kterému rodiče kvůli starcově smůle zakázali s ním dále vyplouvat.\\[10pt] 
\noindent\textbf{Kolize:} Pětaosmdesátý den ráno vyplouvá Santiago na své malé loďce mnohem dál na širé moře než obvykle, aby svou smůlu konečně prolomil. Kolem poledne se mu na návnadu chytí obrovský marlín (mečoun), který je větší než člun, a začne loďku táhnout dál na otevřený oceán.\\[10pt]
\noindent\textbf{Krize:} Probíhá dlouhý, napínavý a vyčerpávající souboj starce s rybou, který trvá tři dny a dvě noci. Santiago je na pokraji fyzických sil, trpí bolestmi, křečemi a únavou. Třetího dne začne ryba kroužit, stařec ji z posledních sil přitáhne k hladině a harpunuje ji. Mrtvého marlína pevně přiváže k boku loďky.\\[10pt]
\noindent\textbf{Peripetie (zvrat):} Zasloužená radost z životního úlovku se náhle mění v zoufalství. Krev mrtvého marlína přiláká žraloky. Santiago s nimi svádí krutý a marný boj. Brání úlovek harpunou, nožem přivázaným k pádlu i kyjem, ale přesila dravců je obrovská a žraloci z nádherné ryby postupně ožerou všechno maso.\\[10pt]
\noindent\textbf{Rozuzlení} Stařec se vrací do přístavu uprostřed noci, naprosto zničený a pouze s obrovskou ohlodanou kostrou. Padne vyčerpáním do postele. Ráno však víc než pětimetrová kostra vzbudí úžas vesničanů, čímž Santiago získává zpět jejich ztracený respekt. Manolin mu v slzách slibuje, že od nynějška budou rybařit opět spolu. Závěr ukazuje, že i když člověk prohraje materiálně, jeho duch nemusí být poražen.\\[10pt]
\mybf{Doba a místo:} na přelomu 40. a 50. let 20. století, Kuba\\[10pt]
\mybf{Vypravěč:} Er-forma\\[10pt]
\mybf{Jazyk:} Jednoduchý, použití šapnělštiny, ale špatného rodu, La mar vs. el mar, jak Santiago mluví o moři. Používá španělský ženský rod (la mar).

\subsection*{Postavy}
\begin{itemize}[itemsep=0.3em]
	\item \textbf{Santiago:} Starý kubánský rybář, ztělesnění lidské nezdolnosti, vytrvalosti a důstojnosti. Je fyzicky poznamenaný těžkým životem, má ruce plné hlubokých jizev, ale jeho oči zůstávají veselé a "neporažené". Přírodu se nesnaží ovládnout, ale cítí se být její součástí.
	
	\item \textbf{Manolin:} Mladý chlapec a Santiagův oddaný učeň. Zastupuje mládí, naději, lásku a soucit. Je se starcem spojen silným mezigeneračním poutem – Santiago mu předává své zkušenosti, chlapec mu na oplátku dává sílu do života a stará se o něj.
	
	\item \textbf{Marlín:} Ačkoliv jde o zvíře, funguje v novele jako plnohodnotná postava. Je to rovnocenný a důstojný soupeř. Symbolizuje krásu přírody, životní výzvu a cíl, za kterým si člověk musí stát.
	
	\item \textbf{Žraloci:} Fungují jako antagonisté příběhu. Představují slepou, destruktivní a krutou sílu přírody (či osudu), která člověka bez lítosti obírá o plody jeho těžce vydřené práce.
\end{itemize}

\subsection*{Autor díla}
\textbf{Eric Arthur Blair (1899--1961)}
Nositel Nobelovy ceny a vůdčí hlas takzvané Ztracené generace, žil život stejně dramatický jako své romány. Byl to nezkrotný dobrodruh, válečný reportér, vášnivý lovec i hlubokomořský rybář, který své zkušenosti z afrických safari, kubánských vod i evropských zákopů přetavil do ikonických děl světové literatury. Zcela změnil moderní prózu svou proslulou „teorií ledovce“ – mimořádně úsporným a strohým stylem, pod jehož na první pohled jednoduchým povrchem se skrývá obrovská tíha nevyřčených emocí. Ačkoliv se však navenek prezentoval jako nezničitelný a tvrdý chlap, v soukromí sváděl vyčerpávající boje s těžkými depresemi, válečnými traumaty a alkoholismem, kterým v roce 1961 nakonec podlehl a svůj život ukončil vlastní rukou.

\textbf{Další díla:} \textit{Sbohem, armádo}, \textit{Komu zvoní hrana}, \textit{Fiesta}.

{\color{subsectioncolor}\hrulefill}
\subsection*{\color{subsectioncolor}Děj}
Novela Ernesta Hemingwaye Stařec a moře vypráví příběh starého kubánského rybáře Santiaga, který už osmdesát čtyři dní neulovil jedinou rybu. Vesničané ho kvůli tomu považují za smolaře a jeho jedinému příteli, chlapci Manolinovi, rodiče zakázali se starcem dál vyplouvat. Santiago se ale nevzdává a pětaosmdesátý den brzy ráno zamíří na své malé loďce mnohem dál na širé moře, než si troufnou ostatní. Věří, že v hlubokých vodách svou vleklou smůlu konečně prolomí.

Kolem poledne pocítí na šňůře silný tah a zjistí, že chytil obrovského marlína (mečouna). Ryba je tak masivní a silná, že ji stařec nedokáže vytáhnout na palubu. Místo toho zvíře začne táhnout loďku daleko za sebou na otevřený oceán. Následuje napínavý a vyčerpávající souboj, který trvá tři dny a dvě noci. Santiago je na vše úplně sám, musí držet napnutou šňůru vlastníma rukama i zádama, trpí bolestivými řeznými ranami, křečemi a nedostatkem spánku. Během těchto těžkých chvil si k rybě vytvoří zvláštní pouto – obdivuje její sílu, oslovuje ji jako svého bratra a cítí k ní hluboký respekt. Třetího dne začne unavený marlín kroužit kolem loďky. Santiago z posledních sil rybu přitáhne k hladině a probodne jí srdce harpunou. Protože je úlovek o víc než půl metru delší než samotný člun a váží přes půl tuny, musí ho pevně přivázat k boku plavidla.

Cesta domů se však brzy změní v boj o přežití. Krev vytékající z mrtvé ryby přiláká žraloky. Prvního dravce stařec ještě dokáže zabít harpunou, ale bohužel o ni při zápasu přijde. Když dorazí další vlny žraloků, brání Santiago svůj životní úlovek s naprostým zoufalstvím – bojuje nožem přivázaným k pádlu, těžkým kyjem i ulomenou pákou kormidla. Přes veškerou jeho obrovskou statečnost je přesila příliš velká a žraloci postupně ožerou z marlína úplně všechno maso.

Santiago se vrací do přístavu uprostřed noci, naprosto vyčerpaný a zlomený. K lodi má přivázanou už jen gigantickou bílou kostru s hlavou a mečem. Odnese si stěžeň do své chatrče a padne do hlubokého spánku. Ráno se kolem člunu seběhnou ostatní rybáři a s ohromným úžasem měří přes pět a půl metru dlouhou kostru. Konečně si uvědomí, jak neskutečný výkon stařec podal, a vrací se mu jejich plný respekt. Když ho v chatrči objeví Manolin, dojetím se rozpláče a pevně starci slíbí, že odteď už budou rybařit zase jen spolu, ať si ostatní říkají, co chtějí. Příběh končí tím, že stařec pokojně spí a znovu se mu zdá o lvech na africkém pobřeží z dob jeho mládí, což dokazuje, že jeho duch nebyl i přes materiální ztrátu nikdy poražen.

\newpage

% ============================================================================
\section{Velký Gatsby}
% ============================================================================

\subsection*{Literární klasifikace}
\mybf{Literární druh:} \\[3pt]
\mybf{Literární forma:} \\[3pt]
\mybf{Funkční styl:} \\[3pt]
\mybf{Slohový postup:} \\[10pt]
\mybf{Myšlenka díla:} \\[10pt]
\mybf{Motiv:} \\[10pt]
\mybf{Inspirace:} \\[10pt]
\mybf{Kompozice:} 

\subsection*{Vnitřní stavba děje}
\noindent\textbf{Expozice} 
\noindent\textbf{Kolize} 
\noindent\textbf{Krize} 
\noindent\textbf{Peripetie (zvrat)} 
\noindent\textbf{Katastrofa} \\[10pt]
\mybf{Doba a místo:} \\[10pt]
\mybf{Vypravěč:} \\[10pt]
\mybf{Jazyk:} 

\subsection*{Postavy}
\begin{itemize}[itemsep=0.3em]
	\item \textbf{Napoleon:}
\end{itemize}

\subsection*{Autor díla}
\textbf{Eric Arthur Blair (1903--1950)}

\textbf{Další díla:} \textit{Barmské dny}, \textit{Farářova dcera}, \textit{1984}, \textit{Na dně v Paříži a Londýně}, \textit{Nadechnout se}.

\subsection*{Literárně historický kontext}

{\color{subsectioncolor}\hrulefill}
\subsection*{\color{subsectioncolor}Děj}

\newpage
% ============================================================================
\section{Farma zvířat}
% ============================================================================

\subsection*{Literární klasifikace}
Dílo lze považovat za více žánrů:
\begin{itemize}[itemsep=0.3em]
	\item \textbf{Antiutopický román:} Antiutopie je opak utopie; představuje fiktivní společnost, která se vyvinula špatným směrem a má zásadní nedostatky (např. totalitní forma vlády, omezování osobní svobody). Občané takového světa jsou často přímo utlačováni politickým systémem.
	\item \textbf{Alegorická bajka:} Příběh kratšího rozsahu, v němž zvířata vystupují jako lidé a zobrazují lidské vlastnosti a chování.
	\item \textbf{Alegorická novela:} Delší alegorické vyprávění, které využívá symboliku a skryté významy k vyjádření kritiky společnosti či politiky.
\end{itemize}
\mybf{Literární druh:} Epika\\[3pt]
\mybf{Literární forma:} Próza\\[3pt]
\mybf{Funkční styl:} Umělecký (literární)\\[3pt]
\mybf{Slohový postup:} Vyprávění\\[10pt]
\mybf{Myšlenka díla:} Hlavním tématem knihy je kritika komunismu a Sovětského svazu. Dílo detailně líčí vývoj totalitního režimu, zpočátku je zde vidina poklidného soužití zvířat, společná hymna \textit{Zvířata Anglie} a sedm přikázání, která mají zaručit spravedlivý řád. Postupně však prasata, která se dostávají k moci, zvířata utlačují a jednají jen ve svůj prospěch, bez ohledu na tvrdou práci ostatních. Kniha ukazuje, že i když je původní myšlenka dobrá, může být zničena, pokud se někteří dostanou k moci, své štěstí a majetek si snaží zajistit i smrt ostatních zvířat a pravidla si přizpůsobují podle svého.\\[10pt]
\mybf{Motiv:} Svoboda, moc, útlak, revoluce, zrada, práce, ideály vs. realita, věrnost, propaganda.\\[10pt]
\mybf{Inspirace:} Inspiroval se dobou stalinistického režimu a v knize skrytě naznačuje, jak funguje.\\[10pt]
\mybf{Kompozice:} Chronologická.

\subsection*{Vnitřní stavba děje}
\noindent\textbf{Expozice} -- Seznamujeme se s prostředím a konfliktem. Ve Veroně se dlouhodobě nenávidí rody Monteků a Kapuletů. Poznáváme hlavní postavy. Romeo a Julie se poprvé potkají na plese u Kapuletů a zamilují se.
\noindent\textbf{Kolize} -- Konflikt se začíná rozvíjet. Romeo a Julie zjistí, že patří k nepřátelským rodům. Přesto se rozhodnou tajně vzít.
\noindent\textbf{Krize} -- Situace se vyhrocuje a komplikuje. Tybalt zabije Mercutia. Romeo v pomstě zabije Tybalta a je vyhnán z Verony. Julie má být provdána za Parise.
\noindent\textbf{Peripetie (zvrat)} -- Dochází k zásadnímu obratu v ději. Julie vypije lektvar od mnicha Vavřince, aby předstírala smrt. Plán je, že Romeo ji potom probudí. Romeovi se ale zpráva nedostane.
\noindent\textbf{Katastrofa} -- Tragický konec. Romeo si myslí, že Julie je opravdu mrtvá, a vypije jed. Julie se probudí a spáchá sebevraždu dýkou. Smrt milenců nakonec usmíří oba rody.\\[10pt]
\mybf{Doba a místo:} Anglický venkov v 50. letech 20. století.\\[10pt]
\mybf{Vypravěč:} Vypravěč je anonymní a příběh je psán z pohledu zvířat. Příběh je psán er-formou (3. osoba) a jedná se o vyprávění a popis.\\[10pt]
\mybf{Jazyk:} Autor používá spisovný jazyk, odborné termíny, metafory, archaismy, personifikace, příběh je psán z pohledu zvířat, er-forma, vypravěč je anonymní.

\subsection*{Postavy}
\begin{itemize}[itemsep=0.3em]
	\item \textbf{Napoleon:} Prase, nejchytřejší, ale je to podvodník
	\item \textbf{Pan Jones:} Původní majitel farmy, alkoholik
	\item \textbf{Molina:} Klisna, ráda se parádí, „dáma"
	\item \textbf{Pištík:} Prase, lhář
	\item \textbf{Major:} Inteligentní kanec
	\item \textbf{Boxer:} Kůň, pracovitý a čestný
	\item \textbf{Kuliš:} Prase, chytré
	\item \textbf{Benjamin:} Osel, pasivní intelektuál
	\item \textbf{Krkavec:} Představuje křesťanství
	\item \textbf{Lidé:} Kapitalisté, vykořisťovatelé
\end{itemize}

\subsection*{Autor díla}
\textbf{Eric Arthur Blair (1903--1950)}, známější pod svým pseudonymem \textbf{George Orwell}, který začal používat, aby si nepřipomínal svou minulost. Britský spisovatel, novinář a esejista. Narodil se v Indii, s matkou se poté přestěhoval do Anglie. Jako mladý vystudoval prestižní školu Eton College. Později se jako dobrovolník zapojil do španělské války. Od roku 1930 pracoval jako novinář. Za války pracoval pro BBC. Ačkoliv kritizoval komunistickou diktaturu, sám se považoval za socialistu. Orwell se zaměřoval na sociální témata své doby, která samozřejmě šla ruku v ruce s politickým režimem. V Česku byl po dlouhou dobu na seznamu zakázaných autorů.

\textbf{Další díla:} \textit{Barmské dny}, \textit{Farářova dcera}, \textit{1984}, \textit{Na dně v Paříži a Londýně}, \textit{Nadechnout se}.

\subsection*{Literárně historický kontext}
\begin{itemize}[itemsep=0.3em]
	\item \textbf{Doba:} Kniha vznikla během 2. světové války, v roce 1945
	\item \textbf{Politická situace:} Stalinismus
	\item \textbf{Rok 1945:} U nás -- moci se pomalu ujímají komunisté, nastupuje politický tlak, proto bylo jeho dílo \textit{Farma zvířat} hned po vydání zakázáno
\end{itemize}

{\color{subsectioncolor}\hrulefill}
\subsection*{\color{subsectioncolor}Děj}
Dílo vypráví o farmě, kde zvířata trpěla hladem a nedostatkem péče, protože pan Jones propíjel peníze a velkou část času trávil v hospodě. Zvířata se vzbouří proti lidské nadvládě a chtějí pracovat jen pro sebe. Revoluce je úspěšná, lidé jsou vyhnáni a farma je přejmenována na Zvířecí farmu. Jsou stanovena pravidla, Sedm přikázání, která nesmějí být porušena.

\begin{center}
	\begin{minipage}{0.8\textwidth}
		\begin{enumerate}[itemsep=0.2em]
			\item Každý, kdo chodí po dvou, je nepřítel.
			\item Všechna zvířata jsou si rovna.
			\item Žádné zvíře nebude chodit oblečené.
			\item Žádné zvíře nebude pít alkohol.
			\item Žádné zvíře nezabije jiné zvíře.
			\item Žádné zvíře nebude spát v posteli.
			\item Každý, kdo chodí po čtyřech nohách, nebo má křídla, je přítel.
		\end{enumerate}
	\end{minipage}
\end{center}

Zvířata začínají pracovat znovu, tentokrát s větším nasazením díky pocitu svobody. V čele stojí dvě prasata, Kuliš a Napoleon, která se neustále hádají. Napoleon nakonec vyžene Kuliše a pošle na něj divoké psy, které si sám vychoval. Postupně nechává zabít i ostatní zvířata a zneužívá jejich zapomnětlivost, aby porušoval Sedm přikázání podle svého. I nejpracovitější zvíře, kůň Boxer, je přetažený a nakonec odvezen na jatka.

Napoleon si vychovává další prasata k svému obrazu a postupně vládnou celé farmě. Nakonec začínají spolupracovat zpět s lidmi, farma se vrací k původnímu názvu Panská farma a prasata se chovají jako lidé, popíjejí a hrají karty. Ostatní zvířata je již nedokáží od lidí rozeznat.

\newpage

% ============================================================================
\section{Krysař}
% ============================================================================

\subsection*{Literární klasifikace}
\mybf{Literární druh:} \\[3pt]
\mybf{Literární forma:} \\[3pt]
\mybf{Funkční styl:} \\[3pt]
\mybf{Slohový postup:} \\[10pt]
\mybf{Myšlenka díla:} \\[10pt]
\mybf{Motiv:} \\[10pt]
\mybf{Inspirace:} \\[10pt]
\mybf{Kompozice:} 

\subsection*{Vnitřní stavba děje}
\noindent\textbf{Expozice} 
\noindent\textbf{Kolize} 
\noindent\textbf{Krize} 
\noindent\textbf{Peripetie (zvrat)} 
\noindent\textbf{Katastrofa} \\[10pt]
\mybf{Doba a místo:} \\[10pt]
\mybf{Vypravěč:} \\[10pt]
\mybf{Jazyk:} 

\subsection*{Postavy}
\begin{itemize}[itemsep=0.3em]
	\item \textbf{Napoleon:}
\end{itemize}

\subsection*{Autor díla}
\textbf{Eric Arthur Blair (1903--1950)}

\textbf{Další díla:} \textit{Barmské dny}, \textit{Farářova dcera}, \textit{1984}, \textit{Na dně v Paříži a Londýně}, \textit{Nadechnout se}.

\subsection*{Literárně historický kontext}

{\color{subsectioncolor}\hrulefill}
\subsection*{\color{subsectioncolor}Děj}

\newpage

% ============================================================================
\section{Válka s mloky}
% ============================================================================

\subsection*{Literární klasifikace}
\mybf{Literární druh:} \\[3pt]
\mybf{Literární forma:} \\[3pt]
\mybf{Funkční styl:} \\[3pt]
\mybf{Slohový postup:} \\[10pt]
\mybf{Myšlenka díla:} \\[10pt]
\mybf{Motiv:} \\[10pt]
\mybf{Inspirace:} \\[10pt]
\mybf{Kompozice:} 

\subsection*{Vnitřní stavba děje}
\noindent\textbf{Expozice} 
\noindent\textbf{Kolize} 
\noindent\textbf{Krize} 
\noindent\textbf{Peripetie (zvrat)} 
\noindent\textbf{Katastrofa} \\[10pt]
\mybf{Doba a místo:} \\[10pt]
\mybf{Vypravěč:} \\[10pt]
\mybf{Jazyk:} 

\subsection*{Postavy}
\begin{itemize}[itemsep=0.3em]
	\item \textbf{Napoleon:}
\end{itemize}

\subsection*{Autor díla}
\textbf{Eric Arthur Blair (1903--1950)}

\textbf{Další díla:} \textit{Barmské dny}, \textit{Farářova dcera}, \textit{1984}, \textit{Na dně v Paříži a Londýně}, \textit{Nadechnout se}.

\subsection*{Literárně historický kontext}

{\color{subsectioncolor}\hrulefill}
\subsection*{\color{subsectioncolor}Děj}

\newpage

% ============================================================================
\section{Bílá nemoc}
% ============================================================================

\subsection*{Literární klasifikace}
\mybf{Literární druh:} \\[3pt]
\mybf{Literární forma:} \\[3pt]
\mybf{Funkční styl:} \\[3pt]
\mybf{Slohový postup:} \\[10pt]
\mybf{Myšlenka díla:} \\[10pt]
\mybf{Motiv:} \\[10pt]
\mybf{Inspirace:} \\[10pt]
\mybf{Kompozice:} 

\subsection*{Vnitřní stavba děje}
\noindent\textbf{Expozice} 
\noindent\textbf{Kolize} 
\noindent\textbf{Krize} 
\noindent\textbf{Peripetie (zvrat)} 
\noindent\textbf{Katastrofa} \\[10pt]
\mybf{Doba a místo:} \\[10pt]
\mybf{Vypravěč:} \\[10pt]
\mybf{Jazyk:} 

\subsection*{Postavy}
\begin{itemize}[itemsep=0.3em]
	\item \textbf{Napoleon:}
\end{itemize}

\subsection*{Autor díla}
\textbf{Eric Arthur Blair (1903--1950)}

\textbf{Další díla:} \textit{Barmské dny}, \textit{Farářova dcera}, \textit{1984}, \textit{Na dně v Paříži a Londýně}, \textit{Nadechnout se}.

\subsection*{Literárně historický kontext}

{\color{subsectioncolor}\hrulefill}
\subsection*{\color{subsectioncolor}Děj}

\newpage

% ============================================================================
\section{Osudy dobrého vojáka Švejka za světové války}
% ============================================================================

\subsection*{Literární klasifikace}
\mybf{Literární druh:} Epika\\[3pt]
\mybf{Literární druh:} Román\\[3pt]
\mybf{Literární forma:} Próza\\[10pt]
\mybf{Myšlenka díla:} Kritika války, armády a nesmyslné poslušnosti autorit.\\[10pt]
\mybf{Inspirace:} Povídky z hospod po Česku\\[10pt]
\mybf{Kompozice:} Chronologická, skládá se ze 4 dílů a ty jsou děleny na kapitoly. Původně mělo být 6 dílů, ale kvůli předčasné smrti napsal Hašek pouze 3, ten 4. dopsal Karel Vaněk\\[10pt]
\mybf{Doba a místo:} První světová válka, Čechy\\[10pt]
\mybf{Vypravěč:} Er-forma\\[10pt]
\mybf{Jazyk:} kombinace spisovné, obecné češtiny a slangu

\subsection*{Postavy}
\begin{itemize}[itemsep=0.3em]
	\item \textbf{Josef Švejk:} „Úřední blb“ a prodejce psů. Svou absurdní horlivostí a doslovným    plněním rozkazů rozkládá rakouskou armádu zevnitř.
	\item \textbf{Nadporučík Lukáš:} Švejkův nešťastný nadřízený. Milovník žen, který marně bojuje se  Švejkovou „pomocí“, jež ho vždy přivede do průšvihu.
	\item \textbf{Polní kurát Katz:} Opilec a hazardní hráč v kněžském rouchu. V Boha nevěří a Švejka  prohraje v kartách. Symbol úpadku církve.
	\item \textbf{Hostinský Palivec:} Majitel hospody „U Kalicha“. Odsouzen na 10 let za poznámku o  mouchách na obraze císaře pána.
	\item \textbf{Agent Bretschneider:} Policejní udavač, který v hospodách provokuje lidi k protistátním řečem. Doplatí na vlastní slídilství.
	\item \textbf{Doktor Grünstein:} Bezcitný vojenský lékař. Každého vojáka považuje za simulanta a  „léčí“ ho klystýrem a hladovkou.
\end{itemize}

\subsection*{Autor díla}
\textbf{Jaroslav Hašek (1883--1923)} -- český spisovatel, novinář, publicista a humorista
Patří mezi nejvýznamnější české satiriky a jeho tvorba vychází z hospodského prostředí a lidového humoru.
Narození - 30.4. 1883, Praha-Nové Město
Úmrtí - 3.1. 1923, Lipnice nad Sázavou
byl rebelem, nedostudoval gymnázium, raději se toulal po Česku a pil v hospodách, kde psal své příběhy
V roce 1915 narukoval do rakousko-uherské armády k 91. pluku
v Českých Budějovicích, byl zajmut Rusy, vstoupil do československých
legií, poté přešel k bolševikům, kde působil jako politický komisař a 
novinář

{\color{subsectioncolor}\hrulefill}
\subsection*{\color{subsectioncolor}Děj}
Příběh začíná v Praze krátce poté, co se Josef Švejk dozvídá o atentátu na arcivévodu Ferdinanda v Sarajevu. Švejk, který se živí prodejem psů s padělanými rodokmeny, přijímá zprávu s klidem a vyráží do hostince U Kalicha. Tam se u piva setkává s policejním udavačem Bretschneiderem, a kvůli neopatrným řečem o mouchách znečišťujících obraz císaře pána je společně s hostinským Palivcem zatčen pro velezradu. Na policejním ředitelství Švejk vyvádí vyšetřovatele z míry svou absolutní ochotou a podepisováním jakýchkoliv přiznání. Po posouzení komisí soudních lékařů, které Švejk ohromuje lidovými moudrostmi a vlasteneckými výkřiky, je úředně prohlášen za „notorického blba“ a převezen do blázince. Přestože je v ústavu naprosto spokojený, lékaři ho brzy označí za simulanta a vyhodí ho zpět na ulici.  Po návratu domů a dalším krátkém zatčení v Salmově ulici, kde vyšetřování paralyzuje proudem historek, se Švejk vrací k prodeji psů. Když však obdrží povolávací rozkaz, odmítá se nechat zastavit záchvatem revmatismu. Nechá se paní Müllerovou vést Prahou na kolečkovém křesle a s nadšeným křikem „Na Bělehrad!“ vyvolává v ulicích obrovský rozruch. Následně putuje do nemocnice pro simulanty k doktoru Grünsteinovi, kde s úsměvem snáší brutální léčbu klystýrem a hladovkou, čímž fanatického lékaře přivádí k zoufalství. Z nemocnice je převezen na garnizónu, kde si ho díky předstíranému dojetí při kázání vybere za osobního sluhu polní kurát Otto Katz.  Služba u Katze, cynického alkoholika, je plná absurdních situací, od vyhazování věřitelů až po improvizovanou polní mši, kde místo vína použijí rum a oltář nahradí kufrem. Vše vrcholí při neúspěšném pokusu o udělení posledního pomazání, kdy Švejk omylem „zaopatří“ lehkou zlomeninu olejem na šicí stroje. Nakonec Katz prohraje Švejka v kartách s nadporučíkem Lukášem. Lukáš se snaží Švejka marně převychovat, ale Švejkova horlivost vede k neustálým pohromám. Definitivní katastrofa nastává, když Švejk pro svého pána ukradne psa, který patří plukovníkovi Schröderovi. Po tomto skandálu jsou oba za trest převeleni k jednotkám do Českých Budějovic a první díl knihy končí jejich nástupem do vlaku směřujícího na frontu.

\newpage

% ============================================================================
\section{Smrt krýsných srnců}
% ============================================================================

\subsection*{Literární klasifikace}
\mybf{Literární druh:} \\[3pt]
\mybf{Literární forma:} \\[3pt]
\mybf{Funkční styl:} \\[3pt]
\mybf{Slohový postup:} \\[10pt]
\mybf{Myšlenka díla:} \\[10pt]
\mybf{Motiv:} \\[10pt]
\mybf{Inspirace:} \\[10pt]
\mybf{Kompozice:} 

\subsection*{Vnitřní stavba děje}
\noindent\textbf{Expozice} 
\noindent\textbf{Kolize} 
\noindent\textbf{Krize} 
\noindent\textbf{Peripetie (zvrat)} 
\noindent\textbf{Katastrofa} \\[10pt]
\mybf{Doba a místo:} \\[10pt]
\mybf{Vypravěč:} \\[10pt]
\mybf{Jazyk:} 

\subsection*{Postavy}
\begin{itemize}[itemsep=0.3em]
	\item \textbf{Napoleon:}
\end{itemize}

\subsection*{Autor díla}
\textbf{Eric Arthur Blair (1903--1950)}

\textbf{Další díla:} \textit{Barmské dny}, \textit{Farářova dcera}, \textit{1984}, \textit{Na dně v Paříži a Londýně}, \textit{Nadechnout se}.

\subsection*{Literárně historický kontext}

{\color{subsectioncolor}\hrulefill}
\subsection*{\color{subsectioncolor}Děj}

\newpage

% ============================================================================
\section{Prima sezona}
% ============================================================================

\subsection*{Literární klasifikace}
\mybf{Literární druh:} \\[3pt]
\mybf{Literární forma:} \\[3pt]
\mybf{Funkční styl:} \\[3pt]
\mybf{Slohový postup:} \\[10pt]
\mybf{Myšlenka díla:} \\[10pt]
\mybf{Motiv:} \\[10pt]
\mybf{Inspirace:} \\[10pt]
\mybf{Kompozice:} 

\subsection*{Vnitřní stavba děje}
\noindent\textbf{Expozice} 
\noindent\textbf{Kolize} 
\noindent\textbf{Krize} 
\noindent\textbf{Peripetie (zvrat)} 
\noindent\textbf{Katastrofa} \\[10pt]
\mybf{Doba a místo:} \\[10pt]
\mybf{Vypravěč:} \\[10pt]
\mybf{Jazyk:} 

\subsection*{Postavy}
\begin{itemize}[itemsep=0.3em]
	\item \textbf{Napoleon:}
\end{itemize}

\subsection*{Autor díla}
\textbf{Eric Arthur Blair (1903--1950)}

\textbf{Další díla:} \textit{Barmské dny}, \textit{Farářova dcera}, \textit{1984}, \textit{Na dně v Paříži a Londýně}, \textit{Nadechnout se}.

\subsection*{Literárně historický kontext}

{\color{subsectioncolor}\hrulefill}
\subsection*{\color{subsectioncolor}Děj}

\newpage

\end{document}